\documentclass[../../main.tex]{subfiles}
\begin{document}

\begin{flushright}
\footnotesize\textit{【自動文字起こし・要確認】}
\end{flushright}

{\bf [解]}
$f_n(x)=x^n\sin^nx$ とおく．さらに，$P(x)=f_1(x)$ とする．はじめに(2)を証明する．

以下，$P(x)=P,\ \sin x=S,\ \cos x=C$ などと書く．
\begin{align}
\begin{split}
S_{n+2}+S_n&>2S_{n+1} \\
&\iff \int_0^{\pi/2}P^n(P^2+1)dx>2\int_0^{\pi/2}P^{n+1}dx \\
&\iff \int_0^{\pi/2}P^n(P-1)^2dx>0 \quad \cdots \text{①}
\end{split}
\end{align}
であり，$[0,\pi/2]$ では $0\le P,\ 0\le(P-1)^2$ だから，
\begin{align}
0\le P^n(P-1)^2
\end{align}
$[0,\pi/2]$ で積分し，常には等号が成立しないことから，
\begin{align}
0<\int_0^{\pi/2}P^n(P-1)^2dx \quad \cdots \text{②}
\end{align}
①，②から(2)が示された．以下(1)，(3)を示す．

ここで $S_1,\ S_2$ を計算する．
\begin{align}
\begin{split}
S_1&=\int_0^{\pi/2}Pdx=\left[-xC+S\right]_0^{\pi/2}=1 \\
S_2&=\int_0^{\pi/2}f_2(x)dx=\int_0^{\pi/2}x^2\cdot\dfrac{1-\cos2x}{2}dx \\
&=\dfrac{1}{2}\left[\dfrac{1}{3}x^3\right]_0^{\pi/2}-\dfrac{1}{2}\left[\dfrac{1}{2}x^2\sin2x+\dfrac{1}{4}\cdot2x\cos2x-\dfrac{1}{8}\cdot2\sin2x\right]_0^{\pi/2} \\
&=\dfrac{1}{6}\left(\dfrac{\pi}{2}\right)^3-\dfrac{1}{2}\cdot\dfrac{1}{4}\cdot2\cdot\left(-\dfrac{\pi}{2}\right)=\dfrac{\pi^3}{48}+\dfrac{\pi}{8}
\end{split}
\end{align}
であり，$\pi>3.14$ から
\begin{align}
\begin{split}
S_2&>\dfrac{\pi}{8}\left(\dfrac{3.14^2}{6}+1\right)>\dfrac{2.64}{8}\cdot3.14=\dfrac{8.2896}{8}>1=S_1 \\
&\therefore\ S_2>S_1 \quad \text{(1)}
\end{split}
\end{align}
となる．したがって，(2)の不等式をくり返し用いて，
\begin{align}
S_{n+1}-S_n>S_n-S_{n-1}>\cdots>S_2-S_1>0 \quad\therefore\ S_{n+1}>S_n
\end{align}
だから，$\{S_n\}$ は単調増加で，$m<n$ ならば $S_m<S_n$ となる． \text{(3)}


\end{document}
